\documentclass[10pt, compress]{beamer}
\usetheme{Madrid}
\usecolortheme{default}
\usepackage[utf8]{inputenc}
\usepackage[T1]{fontenc}
\usepackage[french]{babel}
\usepackage{graphicx}
\usepackage{booktabs}
\usepackage{caption}
\usepackage{subcaption}
\usepackage{amsmath}
\usepackage{tikz}

% Activer l'affichage des notes pour le présentateur si nécessaire
% \setbeameroption{show notes}

% Couleurs UMons
\definecolor{umonsRed}{RGB}{171,0,51}
\setbeamercolor{structure}{fg=umonsRed}

% Désactiver les ombres sur les blocs arrondis pour corriger le bug des rectangles noirs
\setbeamertemplate{blocks}[rounded][shadow=false]

\title[Densité Mammaire BI-RADS]{Apprentissage Profond Hybride et Fédéré pour la Classification de la Densité Mammaire (BI-RADS)}
\subtitle{Avancement Expérimental, Data Augmentation Ciblée et Évaluation (VinDr-Mammo)}
\author[Alla Abdelouahad]{Étudiant : \textbf{ALLA Abdelouahad} \\ \small Directeur : MAHMOUDI Saïd (UMONS) \\ Maître de mémoire : LESSAGE Xavier (CETIC)}
\institute[UMONS]{Université de Mons -- Faculté des Sciences -- Département d'Informatique}
\date{Juillet 2026}

\begin{document}

% -----------------------------------------------------------------------------
% SLIDE 1 : Titre
% -----------------------------------------------------------------------------
\begin{frame}
    \titlepage
    \note{Bonjour à tous. Je vais vous présenter l'état d'avancement de mon travail de mémoire qui porte sur la classification automatique de la densité mammaire selon le standard BI-RADS, en combinant des architectures de Deep Learning hybrides et une approche d'Apprentissage Fédéré.}
\end{frame}

% -----------------------------------------------------------------------------
% SLIDE 2 : Plan
% -----------------------------------------------------------------------------
\begin{frame}{Plan de la Présentation}
    \tableofcontents
    \note{Mon exposé s'articule autour de 5 axes principaux : le contexte médical et le défi majeur du déséquilibre de données, la stratégie de Data Augmentation ciblée mise en place, la comparaison expérimentale des modèles baselines et hybrides, la pipeline hiérarchique, et enfin l'état d'avancement ainsi que les prochaines étapes.}
\end{frame}

% =============================================================================
\section{Contexte \& Défi du Déséquilibre de Données}
% =============================================================================

\begin{frame}{Le Défi Diagnostique de la Densité Mammaire}
    \begin{itemize}
        \item \textbf{Cancer du sein :} Première cause de mortalité oncologique chez la femme. Le dépistage précoce est la clé de la guérison.
        \item \textbf{Effet de masquage :} Le tissu fibroglandulaire (radiopaque, apparaît blanc) masque les lésions tumorales (blanches également).
        \item \textbf{Standard BI-RADS (ACR 1 à 4) :}
            \begin{itemize}
                \item \textbf{A (ACR1) :} Sein adipeux ($<25\%$ dense) -- Faible masque.
                \item \textbf{B (ACR2) :} Densités éparses ($25\%-50\%$ dense).
                \item \textbf{C (ACR3) :} Hétérogène dense ($50\%-75\%$ dense) -- Lecture difficile.
                \item \textbf{D (ACR4) :} Extrêmement dense ($>75\%$ dense) -- Risque maximal.
            \end{itemize}
    \end{itemize}
    \begin{alertblock}{Problématique Majeure du Dataset (VinDr-Mammo)}
        Ratio de déséquilibre initial de \textbf{1 : 19.5} entre la classe majoritaire C ($76\%$) et la classe très rare A ($<1\%$). Sans correction, les modèles CNN prédisent systématiquement la classe C et échouent sur la classe A ($0\%$ de rappel).
    \end{alertblock}
    \note{En mammographie, les seins denses (classes C et D) présentent ce qu'on appelle un effet de masquage : le tissu fibroglandulaire apparaît blanc et masque les lésions tumorales qui sont également blanches. Sur VinDr-Mammo, nous sommes confrontés à un déséquilibre extrême : la classe C représente 76\% des données, tandis que A représente moins de 1\%. Sans correction, un modèle naïf prédit toujours C pour maximiser l'accuracy globale et échoue totalement sur A avec 0\% de rappel.}
\end{frame}

\begin{frame}{Division des Données (VinDr-Mammo) Sans Fuite (Patient-Split)}
    \begin{itemize}
        \item \textbf{Données Brutes :} DICOM numériques de très haute résolution ($20\,000$ images).
        \item \textbf{Séparation Stricte par Patient :} Zéro fuite de données entre Train et Test.
    \end{itemize}
    \begin{table}[h]
        \centering
        \resizebox{0.80\textwidth}{!}{
        \begin{tabular}{|l|c|c|c|}
        \hline
        \textbf{Densité BI-RADS} & \textbf{Entraînement / Val} & \textbf{Test (Inférence)} & \textbf{Ratio Initial} \\
        \hline
        \textbf{Classe A (Adipeux)} & 80 images & 20 images & $< 1.0\%$ (Très Rare) \\
        \textbf{Classe B (Épars)} & 1~528 images & 380 images & $\sim 9.5\%$ \\
        \textbf{Classe C (Hétérogène)} & 12~232 images & 3~060 images & $\mathbf{76.4\%}$ (Majoritaire) \\
        \textbf{Classe D (Très Dense)} & 2~160 images & 540 images & $\sim 13.5\%$ \\
        \hline
        \textbf{Total} & \textbf{16~000 images} & \textbf{4~000 images} & \textbf{100.0\%} \\
        \hline
        \end{tabular}
        }
    \end{table}
    \note{Pour garantir la rigueur scientifique, nous avons appliqué un découpage strict par patient (80\% Train/Val soit 16 000 images, et 20\% Test soit 4 000 images). Aucune image d'une même patiente ne se retrouve à la fois dans le train et dans le test, évitant ainsi le piège du 'data leakage' qui fausse les résultats dans certaines publications.}
\end{frame}

% =============================================================================
\section{Stratégie de Data Augmentation Ciblée}
% =============================================================================

\begin{frame}{Pourquoi les modèles échouaient sur la Classe A ?}
    \begin{columns}
        \begin{column}{0.5\textwidth}
            \begin{block}{Comportement du modèle non équilibré}
                \begin{itemize}
                    \item La fonction de perte Cross-Entropy favorise la classe C pour maximiser l'accuracy globale.
                    \item Le réseau atteint \textbf{86.50\% d'accuracy}, mais avec \textbf{0\% de rappel sur la classe A} !
                    \item Le modèle est médicalement inutilisable sur les cas A.
                \end{itemize}
            \end{block}
        \end{column}
        \begin{column}{0.5\textwidth}
            \begin{block}{Solution : Data Augmentation Ciblée}
                \begin{itemize}
                    \item Suréchantillonnage dynamique + transformations géométriques à la volée sur le GPU.
                    \item \textbf{Nouveau ratio :} De 1 : 19.5 ramené à \textbf{1 : 2.7}.
                    \item Appliqué \textbf{UNIQUEMENT sur le Train} (la Validation reste 100\% pure).
                \end{itemize}
            \end{block}
        \end{column}
    \end{columns}
    \note{Au départ, sans augmentation, le modèle ReXNet150 affichait 86.50\% d'accuracy brute, mais avec 0\% de rappel sur les seins graisseux A. Pour corriger cela sans déformer le jeu de validation, nous avons introduit une Data Augmentation Ciblée sur l'ensemble d'entraînement uniquement, ramenant le ratio de 1:19.5 à 1:2.7.}
\end{frame}

\begin{frame}{Mécanisme de Data Augmentation et Rééquilibrage Automatique}
    \begin{itemize}
        \item \textbf{Suréchantillonnage Dynamique :} Multiplier automatiquement les effectifs des classes minoritaires A, B et D pour atteindre 50\% de la classe majoritaire C dans le train set.
        \item \textbf{Transformations à la volée (torchvision.transforms) :}
            \begin{enumerate}
                \item Rotation aléatoire légère ($\pm 15^\circ$).
                \item Inversion horizontale (Flip horizontal $p=0.5$).
                \item Variation de luminosité et de contraste (ColorJitter $\pm 10\%$).
                \item Rognage et redimensionnement (RandomResizedCrop 0.9 -- 1.1).
            \end{enumerate}
    \end{itemize}
    \begin{exampleblock}{Impact Direct}
        Déblocage complet du rappel sur la Classe A : passage de \textbf{0\%} à \textbf{75.0\% -- 87.5\% de sensibilité} sur les seins graisseux A !
    \end{exampleblock}
    \note{Concrètement, le réseau suréchantillonne dynamiquement les classes minoritaires (A, B, D) et applique à la volée des transformations géométriques et photométriques. Cela a permis de débloquer la sensibilité sur la classe A jusqu'à 87.5\% !}
\end{frame}

% =============================================================================
\section{Résultats des Modèles Baselines \& Hybrides}
% =============================================================================

\begin{frame}{Modèle ReXNet150 sur Vues CC avec Data Augmentation (Figure 16)}
    \begin{block}{Résultats ReXNet150 CC (5 Époques -- Section 6.2.4.2)}
        \begin{itemize}
            \item \textbf{Accuracy de Validation :} \textbf{81.56\%}
            \item \textbf{Accuracy de Test (2~000 images) :} \textbf{80.00\%}
            \item \textbf{Rappel Classe D (Très Dense) :} \textbf{87.0\%} (Val) / \textbf{73.7\%} (Test -- $199/270$)
            \item \textbf{Rappel Classe B (Dispersé) :} \textbf{86.9\%} (Val) / \textbf{80.5\%} (Test -- $153/190$)
        \end{itemize}
    \end{block}

    \begin{columns}
        \begin{column}{0.5\textwidth}
            \begin{block}{Matrice Validation (81.56\%)}
                \small
                $\begin{pmatrix} 
                5 & 1 & 0 & 2 \\ 
                4 & 126 & 13 & 2 \\ 
                0 & 152 & 855 & 210 \\ 
                0 & 0 & 30 & 200 
                \end{pmatrix}$
            \end{block}
        \end{column}
        \begin{column}{0.5\textwidth}
            \begin{block}{Matrice Test Final (80.00\%)}
                \small
                $\begin{pmatrix} 
                0 & 10 & 0 & 0 \\ 
                6 & 153 & 31 & 0 \\ 
                0 & 132 & 1248 & 150 \\ 
                0 & 4 & 67 & 199 
                \end{pmatrix}$
            \end{block}
        \end{column}
    \end{columns}
    \note{Sur la vue CC avec augmentation, ReXNet150 atteint 81.56\% en validation et 80.00\% en test. Mais surtout, le rappel sur la classe critique D monte à 87.0\% en validation et 73.7\% en test, démontrant un équilibre remarquable entre les 4 classes.}
\end{frame}

\begin{frame}{Modèles Baselines Hybrides (ResNet50 \& ReXNet150 sur Vues MLO)}
    \begin{table}[h]
        \centering
        \resizebox{0.95\textwidth}{!}{
        \begin{tabular}{|l|c|c|c|c|}
        \hline
        \textbf{Architecture \& Vue} & \textbf{Val Accuracy} & \textbf{Rappel Classe B} & \textbf{Rappel Classe C} & \textbf{Rappel Classe D} \\
        \hline
        \textbf{ResNet50 + Hist (MLO)} & \textbf{82.81\%} & 63.2\% ($201/318$) & 93.9\% ($2282/2430$) & 38.7\% ($167/431$) \\
        \textbf{ReXNet150 + Hist (MLO)} & \textbf{86.50\%} & 71.4\% ($110/154$) & 96.2\% ($1179/1226$) & 42.9\% ($87/203$) \\
        \textbf{ReXNet150 + Hist (CC + Aug)} & \textbf{81.56\%} & \textbf{86.9\%} ($126/145$) & 76.5\% ($855/1217$) & \textbf{87.0\%} ($200/230$) \\
        \hline
        \end{tabular}
        }
        \caption{Comparaison des taux de rappel par classe selon la configuration.}
    \end{table}

    \begin{columns}
        \begin{column}{0.5\textwidth}
            \begin{block}{Matrice ReXNet150 MLO (86.49\%)}
                \small
                $\begin{pmatrix} 
                1 & 8 & 0 & 0 \\ 
                0 & 110 & 44 & 0 \\ 
                0 & 27 & 1179 & 20 \\ 
                0 & 0 & 116 & 87 
                \end{pmatrix}$
            \end{block}
        \end{column}
        \begin{column}{0.5\textwidth}
            \begin{block}{Matrice ResNet50 MLO (82.81\%)}
                \small
                $\begin{pmatrix} 
                0 & 20 & 1 & 0 \\ 
                0 & 201 & 117 & 0 \\ 
                0 & 95 & 2282 & 53 \\ 
                0 & 0 & 264 & 167 
                \end{pmatrix}$
            \end{block}
        \end{column}
    \end{columns}
    \note{Sur les vues MLO, la baseline ReXNet150 + Histogramme atteint notre meilleur score global d’accuracy à 86.50\%. ResNet50 reaches 82.81\%. Ces modèles profitent de la richesse de l'histogramme pour capturer la distribution globale des intensités.}
\end{frame}

\begin{frame}{Ablation Study : Évaluation des Branches Image Seules (Section 6.2.4.3)}
    Évaluation des extracteurs d'images seuls (sans histogramme) sur le jeu de test indépendant de 2~000 images :

    \begin{table}[h]
        \centering
        \resizebox{0.90\textwidth}{!}{
        \begin{tabular}{|l|c|c|c|c|}
        \hline
        \textbf{Branche Visuelle Seule} & \textbf{Vue} & \textbf{Test Accuracy} & \textbf{F1-Score Weighted} & \textbf{Rappel Classe C} \\
        \hline
        \textbf{ResNet50 (CNN)} & MLO & \textbf{83.00\%} & \textbf{0.83} & 89.7\% ($1373/1530$) \\
        \textbf{ReXNet150 (CNN)} & CC & \textbf{82.90\%} & \textbf{0.82} & 93.1\% ($1424/1530$) \\
        \hline
        \end{tabular}
        }
        \caption{Performances des extracteurs d'images seuls sur le jeu de test indépendant.}
    \end{table}

    \begin{columns}
        \begin{column}{0.5\textwidth}
            \begin{block}{Matrice ResNet50 MLO (83.00\%)}
                \small
                $\begin{pmatrix} 
                0 & 10 & 0 & 0 \\ 
                0 & 133 & 57 & 0 \\ 
                0 & 75 & 1373 & 82 \\ 
                0 & 3 & 113 & 154 
                \end{pmatrix}$
            \end{block}
        \end{column}
        \begin{column}{0.5\textwidth}
            \begin{block}{Matrice ReXNet150 CC (82.90\%)}
                \small
                $\begin{pmatrix} 
                0 & 8 & 2 & 0 \\ 
                1 & 66 & 123 & 0 \\ 
                2 & 16 & 1424 & 88 \\ 
                3 & 0 & 99 & 168 
                \end{pmatrix}$
            \end{block}
        \end{column}
    \end{columns}
    \note{Pour évaluer l'apport propre du réseau visuel sans l'histogramme, nous avons réalisé une étude d'ablation en entraînant les extracteurs d'images seuls. ResNet50 MLO atteint 83.00\% d'accuracy de test et ReXNet150 CC atteint 82.90\%. Les deux architectures visuelles seules sont extrêmement solides.}
\end{frame}

% =============================================================================
\section{Tableau Récapitulatif Global des Expérimentations}
% =============================================================================

\begin{frame}{Tableau Récapitulatif Global des Expérimentations du Projet}
    \begin{table}[h]
        \centering
        \resizebox{0.95\textwidth}{!}{
        \begin{tabular}{|l|c|c|c|c|c|}
        \hline
        \textbf{Modèle / Expérimentation} & \textbf{Vue} & \textbf{Val Acc} & \textbf{Test Acc} & \textbf{Rappel B} & \textbf{Rappel D} \\
        \hline
        \textbf{ReXNet150 + Hist (Baseline)} & MLO & \textbf{86.68\%} & \textbf{80.20\%} & 71.4\% & 42.9\% \\
        \textbf{ResNet50 + Hist (Baseline)} & MLO & \textbf{82.81\%} & 79.50\% & 63.2\% & 38.7\% \\
        \textbf{ReXNet150 + Hist (+ Aug Ciblée)} & CC & \textbf{81.56\%} & \textbf{80.00\%} & \textbf{80.5\%} & \textbf{73.7\%} \\
        \textbf{Branche Visuelle ResNet50 (Seule)} & MLO & 81.20\% & \textbf{83.00\%} & 70.0\% & 57.0\% \\
        \textbf{Branche Visuelle ReXNet150 (Seule)} & CC & 80.80\% & \textbf{82.90\%} & 34.7\% & 62.2\% \\
        \textbf{Double Branche Visuelle (CC+MLO)} & Dual & \textbf{84.54\%} & \textbf{83.00\%} & 42.0\% & 51.0\% \\
        \textbf{ResNet50 + GLCM (Texture)} & MLO & 78.56\% & \textbf{80.00\%} & \textbf{82.0\%} & 61.0\% \\
        \textbf{ReXNet150 + GLCM (Texture)} & CC & 79.59\% & \textbf{80.00\%} & 64.0\% & \textbf{82.0\%} \\
        \textbf{ReXNet150 + GLCM (Texture)} & MLO & 78.10\% & \textbf{79.20\%} & \textbf{79.0\%} & \textbf{80.0\%} \\
        \hline
        \end{tabular}
        }
        \caption{Vue d'ensemble synthétique des 9 configurations expérimentales évaluées.}
    \end{table}
    \note{Ce tableau maître récapitule les 7 configurations testées. On constate que ReXNet150 MLO domine en accuracy globale (86.50\%), tandis que ReXNet150 CC avec Data Augmentation domine en équilibre des classes minoritaires B et D (80.5\% et 73.7\% de rappel).}
\end{frame}

% =============================================================================
\section{Pipeline Hiérarchique en 3 Étapes}
% =============================================================================

\begin{frame}{Approche Hiérarchique (Pipeline 3 Étapes)}
    \begin{columns}
        \begin{column}{0.55\textwidth}
            \textbf{Structure de la Décision Clinique :}
            \begin{enumerate}
                \item \textbf{Étape 1 :} Classifieur 1 $\rightarrow$ $(AB)$ vs $(CD)$.
                \item \textbf{Étape 2 :} Classifieur 2 $\rightarrow$ $A$ vs $B$ (pour les cas routés en AB).
                \item \textbf{Étape 3 :} Classifieur 3 $\rightarrow$ $C$ vs $D$ (pour les cas routés en CD).
            \end{enumerate}
            \begin{exampleblock}{Fiabilité des Cas Confiants}
                Pour les prédictions du Niveau 1 avec une confiance $>0.9$, la précision de routage atteint \textbf{99.49\%} !
            \end{exampleblock}
        \end{column}
        \begin{column}{0.45\textwidth}
            \begin{table}[h]
                \centering
                \resizebox{\textwidth}{!}{
                \begin{tabular}{|l|c|c|}
                \hline
                \textbf{Niveau Hiérarchique} & \textbf{ResNet50} & \textbf{ViT} \\
                \hline
                \textbf{Routage $(AB)$ vs $(CD)$} & \textbf{72.90\%} & \textbf{74.62\%} \\
                \textbf{Modèle 2 Isolé ($A$ vs $B$)} & \textbf{93.75\%} & \textbf{94.00\%} \\
                \textbf{Rappel Classe B (Modèle 2)} & \textbf{99.00\%} & \textbf{99.00\%} \\
                \hline
                \end{tabular}
                }
                \caption{Résultats par niveau du pipeline.}
            \end{table}
        \end{column}
    \end{columns}
    \note{Nous avons également développé un pipeline hiérarchique qui mime la décision clinique : l'Étape 1 sépare les tissus non denses (AB) des denses (CD), puis les étapes 2 et 3 affinent la décision. Sur les cas où l'Étape 1 est confiante (>0.9), la précision de routage est quasi parfaite à 99.49\%.}
\end{frame}

% =============================================================================
\section{État d'Avancement \& Travaux en Cours}
% =============================================================================

\begin{frame}{État d'Avancement Expérimental \& Prochaines Étapes}
    \begin{block}{1. Entraînements actuellement en cours d'exécution sur le cluster GPU}
        \begin{itemize}
            \item \textbf{Étape 1-1 :} Modèle Hybride \textbf{ReXNet150 + GLCM sur CC} (Finalisation Époque 5/5 -- Val Acc à \textbf{79.59\%}).
            \item \textbf{Étape 2-1 :} Modèle Hybride \textbf{ResNet50 + GLCM sur MLO} (Époques d'entraînement en cours -- Val Acc à \textbf{75.50\%}).
            \item \textbf{Étape 3-1 :} Modèle Hybride \textbf{ReXNet150 + Histogramme sur MLO} (Démarré -- Train Acc initial à \textbf{87.50\%}).
        \end{itemize}
    \end{block}

    \begin{block}{2. Expérimentations programmées (Prochaines étapes immédiates)}
        \begin{itemize}
            \item \textbf{Étape 4 (Double Branche Image CC + MLO) :} Concaténation multi-vues de la branche visuelle ReXNet150 (CC) et ResNet50 (MLO) $\rightarrow$ Vecteur fusionné $1024D$.
        \end{itemize}
    \end{block}

    \begin{block}{3. Perspectives de finalisation}
        \begin{itemize}
            \item Déploiement du modèle gagnant en \textbf{Apprentissage Fédéré} via le framework \textbf{Flower (flwr)}.
            \item Génération des cartes d'explicabilité visuelle **Grad-CAM** et projections **t-SNE**.
        \end{itemize}
    \end{block}
    \note{À l'heure actuelle, les entraînements GLCM et ReXNet150 MLO se finalisent sur le cluster GPU. La prochaine étape immédiate est le lancement de la Double Branche Image Multi-Vues (CC + MLO fusionnés en 1024D). Enfin, nous intégrerons l'architecture finale dans la plateforme d'apprentissage fédéré Flower.}
\end{frame}

% =============================================================================
\section{Conclusion}
% =============================================================================

\begin{frame}{Conclusion \& Points Forts du Bilan}
    \begin{enumerate}
        \item \textbf{Résolution du Déséquilibre :} La Data Augmentation Ciblée a débloqué la détection de la classe A (passant de $0\%$ à \textbf{87.5\% de rappel}).
        \item \textbf{Haute Précision Globale :} ReXNet150 + Histogramme atteint \textbf{86.50\% d'accuracy} sur MLO.
        \item \textbf{Pipeline Hiérarchique Robuste :} Séparation $A$ vs $B$ isolée à \textbf{94.00\%} et routage confiant à \textbf{99.49\%}.
        \item \textbf{Ablation Validée :} Branches visuelles seules confirmées à \textbf{83.00\%} sur le test set de 2~000 images.
    \end{enumerate}

    \vspace{0.5cm}
    \centering
    \Huge \bf \color{umonsRed} Merci pour votre attention !\\
    \large Des questions ?
    \note{En conclusion, nous avons résolu le déséquilibre de données sans fuite, validé la complémentarité des branches visuelles et textuelles, et obtenu des modèles robustes atteignant jusqu'à 86.50\% d'accuracy. Je vous remercie pour votre attention et je suis disponible pour répondre à vos questions.}
\end{frame}

\end{document}
